\section{Prox-GNN}
\label{sec:method}

We extend the paragraph-level graph neural network of \citet{ilpcsr} with two modifications informed by the analysis in \S\ref{sec:analysis}. We first describe the base architecture, then the two extensions.

\subsection{Base Architecture: Para-GNN}
\label{sec:base}

We adapt Para-GNN \citep{ilpcsr} for multilingual statute retrieval. For each query-candidate pair, a heterogeneous graph is constructed with four node types: (i) query document nodes, (ii) candidate document nodes, (iii) query paragraph nodes (one per sentence), and (iv) candidate paragraph nodes (one per statute clause). Each paragraph node is connected to its parent document node via a directed edge. Node features are sentence-level embeddings; edge features are rhetorical role embeddings (``NONE'' for statutes). Document node features are the mean of their paragraph embeddings.

A two-layer Edge-Updated Graph Attention Network (EUGAT) processes the graph, updating node representations through learned attention over neighbors and edge features, with residual connections. The final score for a query-candidate pair is:
\begin{equation}
\text{score} = \alpha \cdot \hat{s}_{\text{GNN}} + (1-\alpha) \cdot s_{\text{BM25}}
\label{eq:score}
\end{equation}
where $\hat{s}_{\text{GNN}}$ is the Z-normalized GNN score and $\alpha$ is a blending weight.

We replace the original English-only encoder (all-mpnet-base-v2, 768d) with BGE-M3 \citep{bgem3}, a multilingual encoder (1024d) supporting Indonesian, French, Chinese, and English. All other components (EUGAT layers, training loss, negative sampling) follow \citet{ilpcsr}.

\subsection{Post-Training Alpha Grid Search}
\label{sec:alpha}

In \citet{ilpcsr}, $\alpha$ is learned end-to-end via a feed-forward network. We observe that this learned $\alpha$ converges to approximately 0.97 regardless of dataset, effectively ignoring BM25 (\S\ref{sec:analysis}). However, the optimal $\alpha$ varies from 0.7 to 0.9 depending on the level of lexical overlap in each dataset.

We decouple signal blending from GNN training. After training, we compute the raw GNN scores and BM25 scores for all test query-candidate pairs, then sweep $\alpha \in \{0.0, 0.1, \ldots, 1.0\}$ and select the value that maximizes MRR@10. This takes less than one second as it involves only arithmetic on pre-computed score matrices.

\subsection{Statute Proximity Edges}
\label{sec:proximity}

In the base Para-GNN graph, each statute node is only connected to its own clause nodes. Statute nodes are not connected to each other, even when the corresponding articles are structurally related in the legal code. Our analysis shows that co-relevant statutes are typically close neighbors (median 18 articles apart), while hard negatives are distant (median 297).

We add bidirectional edges between candidate statute nodes whose article numbers are within a radius $N$. Formally, for candidates $s_i$ and $s_j$ with article numbers $a_i$ and $a_j$, we add edges $(s_i \rightarrow s_j)$ and $(s_j \rightarrow s_i)$ if $|a_i - a_j| \leq N$. We set $N=50$ based on the observation that 66.3\% of co-relevant pairs fall within this radius. Edge features use the ``NONE'' embedding. After two EUGAT layers, each statute node's representation incorporates information from its structural neighbors.

This extension is inspired by G-DSR \citep{gdsr}, which models legislation structure as a graph for statute retrieval. The difference is that G-DSR builds a static graph over all statutes, while our proximity edges are within the per-batch candidate graph, allowing the GNN to learn query-specific propagation.

\subsection{Baselines}

We compare against four baselines using the same evaluation protocol.

\noindent\textbf{BM25:} Okapi BM25 with $k_1=1.5$, $b=0.75$. Stemming disabled for Indonesian; jieba segmentation for Chinese (details in App.~\ref{app:implementation}).

\noindent\textbf{Dense (BGE-M3):} Cosine similarity over BGE-M3 dense vectors.

\noindent\textbf{GAR \citep{gar}:} BM25 pool (top 200) re-ranked with a cross-encoder scorer and expanded via corpus graph traversal.

\noindent\textbf{JNLP Stage~1 \citep{jnlp}:} CatBoost classifier on element-wise product features of BGE-M3 query and document embeddings.

\subsection{Evaluation Protocol}

We report MRR@10, Recall@10, and Hit Rate. Each method is evaluated on both KUHPerdata variants (shared relevance judgments) and on BSARD and STARD at their native query levels. All datasets use the same train and test split protocol.
